\chapter{Analyse et spécification des besoins}
\label{ch:analyse-besoins}
\minitoc
\clearpage

\section{Introduction}

L'analyse menée ici se concentre sur les usages réels de PharmacieConnect : la consultation rapide d'une
pharmacie de garde par le citoyen, la mise à jour encadrée des données par l'administration et la validation
des imports CSV. L'enjeu consiste à convertir le besoin général d'accès à l'information pharmaceutique en
règles concrètes, directement exploitables par le backend FastAPI, l'interface React/Vite et l'application
mobile Expo/React Native.

\section{Acteurs}

Les acteurs de PharmacieConnect ont été définis à partir des responsabilités réellement présentes dans
l'application. Le citoyen consulte les données publiques sans intervenir sur leur contenu, tandis que les
comptes d'administration modifient les pharmacies, les gardes et les médicaments à travers des routes
protégées. Cette distinction est indispensable, car une erreur sur un planning de garde ou une localisation
peut rendre l'information inutilisable au moment de la recherche. Le Tableau~\ref{tab:acteurs} récapitule
ces acteurs et leurs responsabilités.

\begin{longtable}{L{4cm}L{10cm}}
    \caption{Acteurs de la plateforme PharmacieConnect\\ \small Source : Auteur}\label{tab:acteurs}\\
    \toprule
    \textbf{Acteur} & \textbf{Rôle dans le système} \\
    \midrule
    \endfirsthead
    \toprule
    \textbf{Acteur} & \textbf{Rôle dans le système} \\
    \midrule
    \endhead
    Administrateur & Supervise les pharmacies, les plannings de garde, le catalogue de médicaments, les imports CSV, les indicateurs, les comptes autorisés et les journaux d'audit. \\
    Assistant ou utilisateur autorisé & Met à jour les pharmacies et les gardes dans un périmètre régional limité, selon les permissions accordées par l'administration. \\
    Utilisateur mobile ou citoyen & Recherche une pharmacie, consulte les gardes par date, utilise la géolocalisation, consulte les médicaments et règle ses préférences d'application. \\
    Système & Valide les données, applique les règles métier, sécurise les accès, calcule les distances, journalise les actions et expose les API aux interfaces. \\
    \bottomrule
\end{longtable}

\section{Besoins fonctionnels détaillés}

Les besoins fonctionnels ont été regroupés selon les modules développés dans le dépôt : authentification,
pharmacies, gardes, médicaments, administration, application mobile et assistant de premiers secours. Ce
découpage correspond aux contrats API utilisés par les clients web et mobile. Il permet également de vérifier
que chaque fonctionnalité décrite dans le rapport possède une traduction concrète dans les routes FastAPI,
les services métier ou les écrans de l'application, comme le détaille le Tableau~\ref{tab:besoins-fonctionnels}.

\begin{longtable}{L{4cm}L{10cm}}
    \caption{Besoins fonctionnels principaux\\ \small Source : Auteur}\label{tab:besoins-fonctionnels}\\
    \toprule
    \textbf{Module} & \textbf{Fonctionnalités attendues} \\
    \midrule
    \endfirsthead
    \toprule
    \textbf{Module} & \textbf{Fonctionnalités attendues} \\
    \midrule
    \endhead
    Authentification & Connexion, inscription, vérification email, rafraîchissement de session, déconnexion, modification du profil et changement de mot de passe. \\
    Recherche de pharmacies & Recherche par texte, consultation paginée, détail d'une pharmacie, recherche par gouvernorat et recherche par proximité à partir des coordonnées GPS. \\
    Pharmacies de garde & Consultation publique par date, création manuelle, modification, suppression et import CSV des plannings côté administration. \\
    Gestion des pharmacies & Ajout, mise à jour, retrait et consultation des fiches, pagination, filtrage régional et import CSV avec retour des erreurs ligne par ligne. \\
    Gestion des médicaments & Import CSV, mise à jour par code PCT, recherche par nom commercial, DCI ou code PCT, consultation du détail et pagination. \\
    Assistant de premiers secours & Réponse extractive aux questions de premiers secours, récupération hybride, score de confiance, avertissement de sécurité et réponse de repli lorsque la confiance est insuffisante. \\
    Administration & Tableau de bord, gestion des assistants, contrôle des autorisations, carte administrative, indicateurs opérationnels et consultation des journaux d'audit. \\
    Application mobile & Accueil, carte, calendrier des gardes, catalogue de médicaments, paramètres, langues, thème sombre, favoris et notifications. \\
    \bottomrule
\end{longtable}

\section{Besoins non fonctionnels}

Les contraintes non fonctionnelles retenues proviennent directement des risques identifiés dans
PharmacieConnect. La sécurité concerne surtout les routes d'administration, les imports CSV et les comptes
assistants. La performance est liée à la pagination des pharmacies et des médicaments. La fiabilité dépend
de la validation des coordonnées GPS, des dates de garde et des colonnes CSV obligatoires. Ces contraintes
orientent donc la conception autant que les fonctionnalités visibles par l'utilisateur.

\begin{itemize}
    \item \textbf{Sécurité} : protéger les routes administratives, hacher les mots de passe, utiliser des jetons JWT,
    vérifier les rôles, limiter les assistants à leur périmètre régional et tracer les actions sensibles.
    \item \textbf{Performance} : paginer les réponses, limiter les tailles d'import, éviter les réponses trop
    volumineuses et utiliser le cache lorsque la configuration le permet.
    \item \textbf{Ergonomie} : proposer une application mobile lisible, une administration claire et une navigation
    adaptée aux opérations fréquentes comme la recherche, l'import et la correction.
    \item \textbf{Maintenabilité} : séparer les routes, services, modèles, schémas, composants et écrans afin de
    faciliter l'évolution du projet.
    \item \textbf{Disponibilité} : rendre les informations de consultation accessibles à travers des endpoints publics
    et prévoir une configuration adaptée au déploiement.
    \item \textbf{Fiabilité des données} : valider les fichiers CSV, contrôler les champs obligatoires, signaler les
    erreurs ligne par ligne et éviter les doublons lorsque des identifiants métier existent.
\end{itemize}

\section{Règles de gestion}

Les règles de gestion ci-dessous correspondent aux contrôles appliqués côté serveur. Elles ne sont pas
laissées aux interfaces clientes, car un appel direct à l'API pourrait contourner une restriction affichée
uniquement dans React ou React Native. Le backend vérifie donc les droits, le périmètre régional, la validité
des coordonnées, la structure des fichiers importés et la journalisation des opérations sensibles avant toute
écriture persistante :

\begin{itemize}
    \item une route d'administration ne peut être appelée que par un compte authentifié et autorisé ;
    \item un assistant ne peut intervenir que dans le périmètre régional qui lui est attribué ;
    \item les mots de passe doivent être hachés avant leur enregistrement ;
    \item les jetons d'accès expirés doivent être renouvelés à travers un jeton de rafraîchissement valide ;
    \item une pharmacie doit posséder au minimum un nom et des informations de localisation exploitables ;
    \item les coordonnées GPS doivent respecter les bornes valides de latitude et de longitude ;
    \item un planning de garde doit être associé à une pharmacie et à une date valide ;
    \item un fichier CSV importé doit respecter les colonnes obligatoires du module concerné ;
    \item les erreurs d'import doivent être retournées avec un message compréhensible ;
    \item les actions sensibles doivent être journalisées dans les logs d'audit avec un acteur, une action, une cible et un statut.
\end{itemize}

\section{Spécification détaillée des cas d'utilisation}

\subsection{Authentification}

Le Tableau~\ref{tab:uc-authentification} décrit ce cas d'utilisation.

\begin{longtable}{L{4cm}L{10cm}}
    \caption{Spécification du cas d'utilisation : authentification\\ \small Source : Auteur}\label{tab:uc-authentification}\\
    \toprule
    \textbf{Élément} & \textbf{Description} \\
    \midrule
    \endfirsthead
    \toprule
    \textbf{Élément} & \textbf{Description} \\
    \midrule
    \endhead
    Acteurs & Administrateur, assistant, utilisateur mobile \\
    Objectif & Donner à l'utilisateur l'accès aux fonctionnalités correspondant à son profil et à ses droits. \\
    Précondition & L'utilisateur possède un compte valide ou suit le processus d'inscription. \\
    Scénario nominal & L'utilisateur saisit ses identifiants depuis le tableau de bord ou l'application mobile. L'API FastAPI vérifie le compte, contrôle son état, enregistre la tentative de connexion, puis émet un jeton d'accès et un jeton de rafraîchissement. Le client conserve la session et n'affiche que les actions autorisées par le rôle retourné. \\
    Postcondition & L'utilisateur est authentifié, une session est créée et seules les fonctionnalités autorisées sont accessibles. \\
    Exceptions & Identifiants invalides, compte désactivé, email non vérifié ou jeton expiré. \\
    \bottomrule
\end{longtable}

\subsection{Recherche d'une pharmacie}

Le Tableau~\ref{tab:uc-recherche} en présente le déroulement.

\begin{longtable}{L{4cm}L{10cm}}
    \caption{Spécification du cas d'utilisation : recherche d'une pharmacie\\ \small Source : Auteur}\label{tab:uc-recherche}\\
    \toprule
    \textbf{Élément} & \textbf{Description} \\
    \midrule
    \endfirsthead
    \toprule
    \textbf{Élément} & \textbf{Description} \\
    \midrule
    \endhead
    Acteur & Utilisateur mobile ou citoyen \\
    Objectif & Retrouver une pharmacie à partir d'un mot-clé, d'un gouvernorat ou de la position GPS. \\
    Précondition & L'application mobile est installée et l'API est accessible. \\
    Scénario nominal & L'utilisateur lance une recherche par nom, adresse ou gouvernorat, ou autorise l'utilisation de sa position GPS. L'application mobile transmet les critères à l'API, qui applique les filtres, vérifie les coordonnées disponibles et calcule la distance lorsque la position est fournie. Les résultats sont ensuite retournés avec les informations utiles à la consultation : nom, adresse, téléphone, gouvernorat et position cartographique. \\
    Postcondition & Les pharmacies correspondantes sont affichées sous forme de liste ou sur la carte. \\
    Exceptions & Position indisponible, recherche vide, absence de résultat ou indisponibilité temporaire de l'API. \\
    \bottomrule
\end{longtable}

\subsection{Consultation des pharmacies de garde}

Le déroulé de ce cas d'utilisation est synthétisé dans le Tableau~\ref{tab:uc-consultation-gardes}.

\begin{longtable}{L{4cm}L{10cm}}
    \caption{Spécification du cas d'utilisation : consultation des gardes\\ \small Source : Auteur}\label{tab:uc-consultation-gardes}\\
    \toprule
    \textbf{Élément} & \textbf{Description} \\
    \midrule
    \endfirsthead
    \toprule
    \textbf{Élément} & \textbf{Description} \\
    \midrule
    \endhead
    Acteur & Utilisateur mobile ou citoyen \\
    Objectif & Consulter les pharmacies de garde pour une date donnée. \\
    Précondition & Des plannings de garde sont enregistrés dans la base. \\
    Scénario nominal & L'utilisateur choisit une date. L'application interroge l'API des gardes et affiche les pharmacies disponibles avec leurs informations utiles. \\
    Postcondition & Les gardes correspondant à la date choisie sont affichées. \\
    Exceptions & Date invalide, absence de planning ou erreur de communication avec l'API. \\
    \bottomrule
\end{longtable}

\subsection{Gestion des pharmacies}

Le Tableau~\ref{tab:uc-gestion-pharmacies} en détaille les étapes.

\begin{longtable}{L{4cm}L{10cm}}
    \caption{Spécification du cas d'utilisation : gestion des pharmacies\\ \small Source : Auteur}\label{tab:uc-gestion-pharmacies}\\
    \toprule
    \textbf{Élément} & \textbf{Description} \\
    \midrule
    \endfirsthead
    \toprule
    \textbf{Élément} & \textbf{Description} \\
    \midrule
    \endhead
    Acteurs & Administrateur, assistant autorisé \\
    Objectif & Ajouter, modifier, consulter ou supprimer les informations des pharmacies. \\
    Précondition & L'acteur est authentifié et possède les droits nécessaires. \\
    Scénario nominal & L'acteur ouvre la page des pharmacies et effectue l'opération souhaitée. Le backend vérifie les droits, valide les données et enregistre la modification. \\
    Postcondition & Les données sont mises à jour et l'action est journalisée lorsqu'elle est sensible. \\
    Exceptions & Données invalides, pharmacie introuvable, droits insuffisants ou périmètre régional non autorisé. \\
    \bottomrule
\end{longtable}

\section{Diagramme de cas d'utilisation}

Le diagramme de cas d'utilisation de PharmacieConnect distingue quatre acteurs. Le \textbf{visiteur},
non authentifié, n'accède qu'aux fonctions publiques : rechercher une pharmacie, consulter les pharmacies
de garde, visualiser la carte et parcourir le catalogue de médicaments. L'\textbf{utilisateur mobile}
dispose en outre d'un compte, ce qui lui ouvre la création de compte, la connexion, la déconnexion, la
gestion de son profil et l'accès à l'assistant de premiers secours. L'\textbf{administrateur} pilote la
plateforme depuis le tableau de bord : gestion des pharmacies, des pharmacies de garde et du catalogue de
médicaments, imports CSV, gestion des assistants et des utilisateurs, et consultation des journaux d'audit.
L'\textbf{assistant régional} partage un sous-ensemble de ces opérations -- pharmacies, gardes et imports
CSV -- mais restreint à son périmètre. Comme le résume la Figure~\ref{fig:use-case-detaille}, les actions
sensibles s'appuient sur une authentification (\og se connecter \fg{}) et une vérification des rôles RBAC,
ce qui matérialise la séparation entre les actions publiques et les actions protégées.

\begin{figure}[H]
    \centering
    \includegraphics[width=0.82\textwidth,keepaspectratio]{usecasediag}
    \caption{Diagramme de cas d'utilisation détaillé\\ \small Source : Auteur}
    \label{fig:use-case-detaille}
\end{figure}

\section{Conclusion}

Au terme de cette analyse, les acteurs, les besoins et les règles de gestion de PharmacieConnect sont
posés. Il en ressort que la plateforme doit concilier consultation rapide, administration encadrée, qualité
des données et traçabilité. Le chapitre suivant traduit ces besoins en architecture, modèle de données,
diagrammes et mécanismes de sécurité.
